% GENERATED FILE. Edit canonical lessons, then run npm run generate:books.
% canonical-source-hash: fnv1a64:3cd487333799f864
% canonical-lessons: LA-C481-convallis, LA-C481-planities, LA-C481-saxosus, LA-C481-confragosus, LA-C481-delabor

\chapter{Enclosed valley, Plain, Rocky, Rough, To slide down}
\label{ch:la-enclosed-valley-plain-rocky-rough-to-slide-down}
\hlchaptermodality{481}

\begin{chapteropening}
\textbf{By the end of this chapter:} \emph{I can describe the shape of the land: an enclosed valley, a plain, rocky and rough ground, and sliding down a slope.}

Complete the last lesson of chapter 481: I can describe the shape of the land: an enclosed valley, a plain, rocky and rough ground, and sliding down a slope.
\end{chapteropening}

\section[convallis, convallis]{convallis, convallis — an enclosed valley, a deep valley}

\label{lesson:LA-C481-convallis}



\begin{quote}
Before the new one: say the Latin for steep, then the Latin for to climb.
\end{quote}

\subsection*{You'll want to know: convallis, convallis}
\textbf{convallis, convallis} — \textquotedblleft{}an enclosed valley, a deep valley\textquotedblright{}.

\begin{quote}
Inter montēs convallis viridis est.
\end{quote}

\begin{quote}
\emph{Between the mountains there is a green valley.}
\end{quote}

Third declension, feminine: \emph{convallis, convallis f.} A valley shut in on all sides.

\begin{cousinweb}
\emph{convallis} is built from \emph{con-} 'together, all round' and \emph{vallis} 'valley' (EtymDB; Lewis \& Short).
\end{cousinweb}

\begin{grammarlens}[title={What you've built}]
One more word for telling, explaining and describing at B1.
\end{grammarlens}

\subsection*{Your turn}
\begin{itemize}
\raggedright
  \item \emph{Say it:} \emph{convallis, convallis}
  \item \emph{Say it:} \emph{convallis, convallis}, once more
  \item \emph{Say it:} a sentence of your own with \emph{convallis, convallis}
  \item \emph{Recall:} say the Latin for steep, then the Latin for to climb, then say \emph{convallis, convallis} again
\end{itemize}

\begin{tcolorbox}[breakable,colback=teal!4,colframe=teal!35!black,title={Before you move on}]
What does convallis, convallis mean? (\textquotedblleft{}An enclosed valley, a deep valley\textquotedblright{}.) Say it once more.
\end{tcolorbox}
\section[plānitiēs, plānitiēī]{plānitiēs, plānitiēī — a plain, level ground}

\label{lesson:LA-C481-planities}



\begin{quote}
Before the new one: say the Latin for to climb, then the Latin for an enclosed valley.
\end{quote}

\subsection*{You'll want to know: plānitiēs, plānitiēī}
\textbf{plānitiēs, plānitiēī} — \textquotedblleft{}a plain, level ground\textquotedblright{}.

\begin{quote}
Post montēs plānitiēs magna est, et via rēcta est.
\end{quote}

\begin{quote}
\emph{Beyond the mountains there is a great plain, and the road is straight.}
\end{quote}

Fifth declension, feminine: \emph{plānitiēs, plānitiēī f.}, from \emph{plānus} 'flat'. English \emph{plain} comes from \emph{plānus}.

\begin{cousinweb}
\emph{plānitiēs} is built on \emph{plānus} 'flat, level' with the suffix \emph{-itiēs} (Lewis \& Short). English \emph{plain} and \emph{plane} come from \emph{plānus}.
\end{cousinweb}

\begin{grammarlens}[title={What you've built}]
One more word for telling, explaining and describing at B1.
\end{grammarlens}

\subsection*{Your turn}
\begin{itemize}
\raggedright
  \item \emph{Say it:} \emph{plānitiēs, plānitiēī}
  \item \emph{Say it:} \emph{plānitiēs, plānitiēī}, once more
  \item \emph{Say it:} a sentence of your own with \emph{plānitiēs, plānitiēī}
  \item \emph{Recall:} say the Latin for to climb, then the Latin for an enclosed valley, then say \emph{plānitiēs, plānitiēī} again
\end{itemize}

\begin{tcolorbox}[breakable,colback=teal!4,colframe=teal!35!black,title={Before you move on}]
What does plānitiēs, plānitiēī mean? (\textquotedblleft{}A plain, level ground\textquotedblright{}.) Say it once more.
\end{tcolorbox}
\section[saxōsus, saxōsa, saxōsum]{saxōsus, saxōsa, saxōsum — rocky, stony}

\label{lesson:LA-C481-saxosus}



\begin{quote}
Before the new one: say the Latin for an enclosed valley, then the Latin for a plain.
\end{quote}

\subsection*{You'll want to know: saxōsus, saxōsa, saxōsum}
\textbf{saxōsus, saxōsa, saxōsum} — \textquotedblleft{}rocky, stony\textquotedblright{}.

\begin{quote}
Ager saxōsus est; agricola eum arāre nōn potest.
\end{quote}

\begin{quote}
\emph{The field is stony; the farmer cannot plough it.}
\end{quote}

An adjective of three endings, from \emph{saxum} 'rock' with \emph{-ōsus} 'full of'.

\begin{cousinweb}
\emph{saxōsus} is built on \emph{saxum} 'rock, stone' with the suffix \emph{-ōsus} 'full of' (Lewis \& Short).
\end{cousinweb}

\begin{grammarlens}[title={What you've built}]
One more word for telling, explaining and describing at B1.
\end{grammarlens}

\subsection*{Your turn}
\begin{itemize}
\raggedright
  \item \emph{Say it:} \emph{saxōsus, saxōsa, saxōsum}
  \item \emph{Say it:} \emph{saxōsus, saxōsa, saxōsum}, once more
  \item \emph{Say it:} a sentence of your own with \emph{saxōsus, saxōsa, saxōsum}
  \item \emph{Recall:} say the Latin for an enclosed valley, then the Latin for a plain, then say \emph{saxōsus, saxōsa, saxōsum} again
\end{itemize}

\begin{tcolorbox}[breakable,colback=teal!4,colframe=teal!35!black,title={Before you move on}]
What does saxōsus, saxōsa, saxōsum mean? (\textquotedblleft{}Rocky, stony\textquotedblright{}.) Say it once more.
\end{tcolorbox}
\section[cōnfragōsus, cōnfragōsa, cōnfragōsum]{cōnfragōsus, cōnfragōsa, cōnfragōsum — rough, broken (of ground)}

\label{lesson:LA-C481-confragosus}



\begin{quote}
Before the new one: say the Latin for a plain, then the Latin for rocky.
\end{quote}

\subsection*{You'll want to know: cōnfragōsus, cōnfragōsa, cōnfragōsum}
\textbf{cōnfragōsus, cōnfragōsa, cōnfragōsum} — \textquotedblleft{}rough, broken (of ground)\textquotedblright{}.

\begin{quote}
Locus cōnfragōsus est, et equī cadunt.
\end{quote}

\begin{quote}
\emph{The ground is rough, and the horses stumble.}
\end{quote}

An adjective of three endings: ground full of breaks and holes. Also of a hard question: knotty.

\begin{cousinweb}
\emph{cōnfragōsus} is built from \emph{con-} and \emph{fragōsus} 'broken, rough' (EtymDB), which comes from \emph{frangō} 'to break' (Lewis \& Short).
\end{cousinweb}

\begin{grammarlens}[title={What you've built}]
One more word for telling, explaining and describing at B1.
\end{grammarlens}

\subsection*{Your turn}
\begin{itemize}
\raggedright
  \item \emph{Say it:} \emph{cōnfragōsus, cōnfragōsa, cōnfragōsum}
  \item \emph{Say it:} \emph{cōnfragōsus, cōnfragōsa, cōnfragōsum}, once more
  \item \emph{Say it:} a sentence of your own with \emph{cōnfragōsus, cōnfragōsa, cōnfragōsum}
  \item \emph{Recall:} say the Latin for a plain, then the Latin for rocky, then say \emph{cōnfragōsus, cōnfragōsa, cōnfragōsum} again
\end{itemize}

\begin{tcolorbox}[breakable,colback=teal!4,colframe=teal!35!black,title={Before you move on}]
What does cōnfragōsus, cōnfragōsa, cōnfragōsum mean? (\textquotedblleft{}Rough, broken (of ground)\textquotedblright{}.) Say it once more.
\end{tcolorbox}
\section[dēlābor, dēlābī]{dēlābor, dēlābī — to slide down, to slip down, to fall}

\label{lesson:LA-C481-delabor}



\begin{quote}
Before the new one: say the Latin for rocky, then the Latin for rough.
\end{quote}

\subsection*{You'll want to know: dēlābor, dēlābī}
\textbf{dēlābor, dēlābī} — \textquotedblleft{}to slide down, to slip down, to fall\textquotedblright{}.

\begin{quote}
Nix dē monte dēlābitur.
\end{quote}

\begin{quote}
\emph{Snow slides down from the mountain.}
\end{quote}

Deponent: \emph{dēlābor, dēlābī, dēlāpsus sum}, with \emph{dē} or \emph{ex}.

\begin{cousinweb}
\emph{dēlābor} is built from \emph{dē-} 'down' and \emph{lābor} 'to slip, to glide' (EtymDB; Lewis \& Short). English \emph{lapse} comes from \emph{lābor}.
\end{cousinweb}

\begin{grammarlens}[title={What you've built}]
One more word for telling, explaining and describing at B1.
\end{grammarlens}

\subsection*{Your turn}
\begin{itemize}
\raggedright
  \item \emph{Say it:} \emph{dēlābor, dēlābī}
  \item \emph{Say it:} \emph{dēlābor, dēlābī}, once more
  \item \emph{Say it:} a sentence of your own with \emph{dēlābor, dēlābī}
  \item \emph{Recall:} say the Latin for rocky, then the Latin for rough, then say \emph{dēlābor, dēlābī} again
\end{itemize}

\begin{tcolorbox}[breakable,colback=teal!4,colframe=teal!35!black,title={Before you move on}]
What does dēlābor, dēlābī mean? (\textquotedblleft{}To slide down, to slip down, to fall\textquotedblright{}.) Say it once more.
\end{tcolorbox}
